\chapter*{Заключение}
\addcontentsline{toc}{chapter}{Заключение}

В результате выполнения выпускной квалификационной работы разработана и экспериментально проверена обучаемая нейросетевая модель для синтеза и оценки мелодий по биологическим данным. Термин ``самообучающаяся модель'', заданный в теме работы, в итоговой реализации трактуется как модель, параметры которой автоматически оптимизируются на данных. Первоначально рассматриваемая генеративно-состязательная схема с отдельным оценщиком была заменена на более устойчивую архитектуру Bio-Conditioned Transformer, поскольку для рассматриваемой задачи важны воспроизводимость обучения, интерпретируемая функция потерь и возможность отдельно проверять качество результата внешними метриками.

В работе изучены существующие подходы к биосонификации и генерации символической музыки. Анализ показал, что прямое отображение нуклеотидов или аминокислот в ноты обычно не обеспечивает музыкальной связности, а генеративные методы без структурного разделения плохо контролируют гармонию и мелодию. Поэтому в разработанной системе биологическая последовательность не преобразуется в музыку напрямую: сначала из неё извлекаются признаки, затем эти признаки используются как управляющий сигнал для генерации гармонии и мелодии.

Основным результатом является программная система BioSonification с иерархической архитектурой Bio->Harmony->Melody. В системе реализованы биологический энкодер для DNA, RNA и protein последовательностей, модуль структурирования MIDI-корпуса, pairing биологических и музыкальных сегментов, модели генерации гармонии и мелодии, MIDI-renderer и веб-интерфейс для пользовательской генерации. Исходный код и структура проекта документированы; описание программы, руководство пользователя и руководство разработчика приведены в приложении А.

Система обучена на корпусе POP909 и геномных данных RefSeq. Актуальная модель использовала 3746 биологических фрагментов, 25940 музыкальных сегментов и 15925 обучающих пар. Обучение выполнено на GPU NVIDIA RTX 2060 6GB с использованием mixed precision. Модель показала validation loss 0.1454 для гармонии и 0.1566 для мелодии, что лучше показателей ранней дипломной версии модели.

Экспериментальная оценка подтвердила, что разработанная система создаёт валидные двухдорожечные MIDI-файлы. Сравнение актуальной модели с предыдущей дипломной моделью и random baseline показало рост гармонической согласованности мелодии: chord-tone ratio составил 0.6710 против 0.6073 у предыдущей модели и 0.4043 у случайного baseline. При этом абсолютные показатели, такие как число нот или энтропия pitch class, не следует трактовать как универсальную оценку художественного качества: они отражают технические свойства MIDI-фрагментов и зависят от режима генерации.

Поставленная цель достигнута: разработана и проверена обучаемая нейросетевая модель, которая использует биологические признаки как управляющий сигнал для структурированной музыкальной генерации. Решены задачи анализа предметной области, подготовки данных, извлечения признаков, проектирования и реализации архитектуры, обучения моделей, генерации MIDI-файлов и экспериментальной оценки качества.

Разработанное решение имеет ограничения. Модель зависит от музыкального стиля корпуса POP909, генерирует символическое MIDI-представление, использует монофоническую мелодию и ограниченный набор аккордовых качеств. Веб-генерация строится из коротких 4-тактовых фрагментов, поэтому она хорошо подходит для воспроизводимого пользовательского сценария, но не формирует полноценную длинную музыкальную форму без последующей аранжировки. Кроме того, результаты работы не доказывают причинную связь между биологией и музыкой и не означают, что музыкальная структура содержится в ДНК буквально. Корректная интерпретация состоит в том, что биологические признаки могут быть использованы как управляющее условие для генеративной музыкальной модели.

Дальнейшее развитие системы может быть связано с расширением музыкального корпуса, использованием предобученных биологических embeddings, увеличением длины связной генерации, добавлением аккомпанемента, динамики и артикуляции, а также с проведением пользовательского или экспертного эксперимента для оценки музыкального восприятия результатов. Отдельным важным направлением является абляционная проверка real bio, shuffled bio и neutral bio, которая позволит точнее отделить вклад биологического управляющего сигнала от музыкального prior модели.

Таким образом, работа показывает, что биосонификация может быть реализована не только как прямое отображение символов в ноты, но и как управляемая нейросетевая генерация. Предложенная архитектура обеспечивает воспроизводимый путь от биологической последовательности к музыкально связному MIDI-фрагменту, а экспериментальная оценка подтверждает улучшение технической согласованности гармонии и мелодии по сравнению со случайным baseline и ранней версией модели.

\vspace{8mm}

Выпускная квалификационная работа выполнена мной самостоятельно и с соблюдением правил профессиональной этики. Все использованные в работе материалы и заимствованные принципиальные положения (концепции) из опубликованной научной литературы и других источников имеют ссылки на них. Я несу ответственность за приведенные данные и сделанные выводы.

Я ознакомлен с программой государственной итоговой аттестации, согласно которой обнаружение плагиата, фальсификации данных и ложного цитирования является основанием для не допуска к защите выпускной квалификационной работы и выставления оценки <<неудовлетворительно>>.

\vspace{18mm}

\noindent
\begin{minipage}[t]{0.42\textwidth}
\centering
\rule{\linewidth}{0.4pt}\\
\textit{ФИО студента}
\end{minipage}
\hfill
\begin{minipage}[t]{0.42\textwidth}
\centering
\rule{\linewidth}{0.4pt}\\
\textit{Подпись студента}
\end{minipage}

\vspace{10mm}

\noindent <<\uline{\hspace{10mm}}>>\uline{\hspace{32mm}}20\uline{\hspace{7mm}}г.
